%==========================================================================
%  TRES0312 Fysikk — Foiler-mal
%
%  MØrk/lys modus: kommenter ut én av de to linjene nedenfor.
%==========================================================================
\newif\ifdarkmode
\darkmodetrue      % ← mørk modus
%\darkmodefalse   % ← lys modus  (fjern % foran og legg % foran linjen over)

\documentclass[aspectratio=169, 12pt]{beamer}

\usepackage{amd-beamer}
\usetikzlibrary{overlay-beamer-styles}
\usetikzlibrary{calc}

%--------------------------------------------------------------------------
%  DOKUMENTINFO
%--------------------------------------------------------------------------
\title{Dynamikk}
\subtitle{TRES0312 --- Krefter i sirkelbevegelse}
\author{Ben David Normann}
\date{\today}

%--------------------------------------------------------------------------
%  SEKSJONSKONTEKST
%--------------------------------------------------------------------------
\setcounter{section}{4}   % Kapittel 4: "Dynamikk"

\begin{document}

%--------------------------------------------------------------------------
%  SLIDE 1 — TITTELSIDE
%--------------------------------------------------------------------------
\begin{frame}[plain]
  \titlepage
  \dekkerdelkap{kapittel 4.4.2}
  \vspace{0.3cm}
  \begin{center}
  {\bfseries\color{repcol} Første time: sentripetalkraft. Etter pausen: dosert veibane.}
  \end{center}
\end{frame}

%--------------------------------------------------------------------------
%  SLIDE 2 — Tre kjappe fra forrige gang (skråplan)
%--------------------------------------------------------------------------
\begin{frame}{}

  \repetisjon{Tre kjappe fra forrige gang}{%
    \begin{enumerate}[1.]
      \item Hva blir størrelsen på normalkraften $\vec{N}$ på en kloss på et skråplan med helningsvinkel $\theta$, uttrykt ved $m$, $\mathrm{g}$ og $\theta$?
      \item Skriv opp fem-stegs-algoritmen for å løse sammensatte kraftoppgaver.
      \item Hva er formelen for kritisk vinkel $\theta_c$, og hva betyr den fysisk?
    \end{enumerate}
  }

\end{frame}

%--------------------------------------------------------------------------
%  SLIDE 3 — Definisjon: Sentripetalkraft (med figur)
%--------------------------------------------------------------------------
\begin{frame}{Sentripetalkraft}
  \begin{columns}[c]
    \begin{column}{0.55\textwidth}
      \defnat{4.10}{Sentripetalkraft}{
      Kraftsummen for sirkelbevegelse med radius $r$ og konstant banefart $v$ er gitt ved
      \[\sum F = m\cdot\frac{v^2}{r},\]
      og er rettet inn mot sentrum av sirkelbevegelsen.
      }
    \end{column}
    \begin{column}{0.42\textwidth}
      \begin{figure}[ht!]
      \centering
      \begin{tikzpicture}[font=\small]
        \draw[dashed, gray!60, line width=0.8pt] (0,0) circle (1.5);
        \fill[black!60] (0,0) circle (2pt);
        \node[font=\scriptsize, below left=2pt] at (0,0) {sentrum};
        \draw[headCol!70!black, line width=1.4pt] (0,0) -- (1.5,0);
        \fill[headCol!40] (1.5,0) circle (0.18);
        \draw[headCol!70!black, line width=1pt] (1.5,0) circle (0.18);
        \node[font=\scriptsize, right=3pt] at (1.68,0) {$m$};
        \draw[-{Stealth[length=10pt,width=6pt]}, line width=2pt, defscol]
          (1.32,0) -- (0.35,0);
        \node[defscol, font=\small\bfseries, above=3pt] at (0.84,0) {$\sum\vec{F}$};
        \draw[-{Stealth[length=10pt,width=6pt]}, line width=2pt, oppsumcol]
          (1.5,0.18) -- (1.5,1.10);
        \node[oppsumcol, font=\small\bfseries, right=3pt] at (1.5,0.64) {$\vec{v}$};
        \draw[-{Stealth[length=5pt,width=4pt]}, gray!70, thin] (0,0) -- (135:1.5);
        \node[font=\scriptsize, gray!70, above right=0pt] at (135:0.75) {$r$};
      \end{tikzpicture}
      \caption{\footnotesize Sett ovenfra: ball i sirkelbane. Kraftsummen $\sum\vec{F}$ er alltid rettet inn mot sentrum.}
      \end{figure}
    \end{column}
  \end{columns}
\end{frame}

%--------------------------------------------------------------------------
%  SLIDE 4 — Aktivitet: Slegge
%--------------------------------------------------------------------------
\begin{frame}{Aktivitet}

  \aktivitet{Slegge}{%
  En sleggekaster svinger en slegge med masse $m = 4.0\,\mathrm{kg}$ rundt seg. Avstanden fra hånda til kula er $r = 1.2\,\mathrm{m}$, og kula holder konstant fart $v = 15\,\mathrm{m/s}$. Vi ser bort fra tyngdekraften.

  Finn kraften i wiren.
  }

\end{frame}

%--------------------------------------------------------------------------
%  SLIDE 5 — Aktivitet: Karusell
%--------------------------------------------------------------------------
\begin{frame}{Aktivitet}

  \aktivitet{Karusell}{%
  Et barn med masse $m = 25\,\mathrm{kg}$ sitter $r = 3.0\,\mathrm{m}$ fra sentrum av en karusell. Karusellen snurrer med konstant fart $v = 4.0\,\mathrm{m/s}$.

  Hva er den resulterende sentripetalkraften på barnet?
  }

\end{frame}

%--------------------------------------------------------------------------
%  SLIDE 6 — Aktivitet: Bil i kurve
%--------------------------------------------------------------------------
\begin{frame}{Aktivitet}

  \aktivitet{Bil i kurve}{%
  En bil med masse $m = 1200\,\mathrm{kg}$ kjører gjennom en horisontal kurve med radius $r = 50\,\mathrm{m}$ og konstant fart $v = 12\,\mathrm{m/s}$.
  \begin{enumerate}[a)]
      \item Finn den nødvendige sentripetalkraften.
      \item Hvilken kraft er det i praksis som gir denne sentripetalkraften på en flat, horisontal vei?
  \end{enumerate}
  }

\end{frame}

%--------------------------------------------------------------------------
%  SLIDE 7 — Aktivitet: Vogn i loop på tivoli
%--------------------------------------------------------------------------
\begin{frame}{Aktivitet}
  \begin{columns}[c]
    \begin{column}{0.58\textwidth}
      \aktivitet{Vogn i loop}{%
      På tivoli går en vogn i en vertikal, helt sirkulær loop. Figuren viser fart i bunnen av loopen (A), halvveis opp loopen (B), på toppen av loopen (C) og halvveis ned igjen (D).
      \begin{enumerate}[a)]
          \item Finn uttrykk for normalkraften i alle tilfeller.
          \item Hva er betingelsen på $N_\textrm{C}$ for at vogna ikke skal falle av på toppen?
      \end{enumerate}
      }
    \end{column}
    \begin{column}{0.40\textwidth}
      \centering
      \begin{tikzpicture}[font=\small, scale=1.1]
        % Loopen (skinnegang)
        \draw[gray!70, line width=1.6pt] (0,0) circle (1.6);
        % Sentrum
        \fill[gray!70] (0,0) circle (1.8pt);
        % Radius r (enveis pil fra sentrum)
        \draw[-{Stealth[length=5pt,width=4pt]}, gray!70, thin] (0,0) -- (135:1.6);
        \node[font=\scriptsize, gray!70, above right=0pt] at (135:0.8) {$r$};
        % Vogna i A (bunn), B (halvveis opp), C (topp), D (halvveis ned)
        \foreach \ang/\pos in {270/A, 0/B, 90/C, 180/D}{
          \begin{scope}[shift={(\ang:1.6)}, rotate=\ang+90]
            % Lokal y-akse peker inn mot sentrum, lokal x-akse langs fartsretningen
            \fill[headCol!40] (-0.26,0.07) rectangle (0.26,0.34);
            \draw[headCol!70!black, line width=0.8pt] (-0.26,0.07) rectangle (0.26,0.34);
            \fill[headCol!70!black] (-0.15,0.05) circle (0.05);
            \fill[headCol!70!black] ( 0.15,0.05) circle (0.05);
            % Fart (tangentielt, på utsiden av loopen)
            \draw[-{Stealth[length=7pt,width=5pt]}, line width=1.6pt, oppsumcol]
              (-0.1,-0.16) -- (0.85,-0.16);
            \node[oppsumcol, font=\small] at (0.45,-0.45) {$v_\text{\pos}$};
            % Posisjonsnavn (inne i loopen)
            \node[font=\small\bfseries] at (0,0.72) {\pos};
          \end{scope}
        }
      \end{tikzpicture}
    \end{column}
  \end{columns}
\end{frame}

%--------------------------------------------------------------------------
%  SLIDE 8 — PAUSE
%--------------------------------------------------------------------------
\begin{frame}
\begin{center}
  \Huge{Pause}\\[8pt]
  {\large Etter pausen: dosert veibane}
\end{center}
\end{frame}

%--------------------------------------------------------------------------
%  SLIDE 9 — Aktivitet: Dosering av veibane (figur etter pause)
%--------------------------------------------------------------------------
\begin{frame}{Aktivitet}
  \begin{columns}[c]
    \begin{column}{0.52\textwidth}
      \aktivitet{Dosering av veibane}{%
      Er det mulig for en bil med konstant fart $v$ å kjøre gjennom en sirkelformet sving med radius $r$ uten friksjon mellom dekk og vei? Bestem nødvendig dosering (helningsvinkel) $\alpha$. Bruk $v = 90\,\mathrm{km/h}$ og $r = 200\,\mathrm{m}$.
      }
    \end{column}
    \begin{column}{0.46\textwidth}
      \pause
      \begin{figure}[ht!]
      \centering
  \begin{tikzpicture}[font=\small, scale=1.4]
      \fill[black!12] (0.0,0.0) -- (3.5,1.29) -- (3.5,0.0) -- cycle;
      \draw[very thick] (0.0,0.0) -- (3.5,1.29);
      \draw[black!30,thin] (0.0,0.0) -- (3.5,0.0);
      \draw[black!30,thin] (3.5,0.0) -- (3.5,1.29);
      \draw[black!55,thin] (0.60,0) arc (0:22:0.60);
      \node[font=\scriptsize,black!70] at (0.74,0.13) {$\alpha$};
      \begin{scope}[rotate=22]
        \fill[headCol!20] (1.50,0.0) rectangle (2.40,0.50);
        \draw[headCol!65!black,thick] (1.50,0.0) rectangle (2.40,0.50);
        \node[font=\scriptsize] at (1.95,0.25) {$m$};
      \end{scope}
      \draw[-{Stealth[length=11pt,width=7pt]}, line width=2.5pt, defscol]
        (1.808,0.731) -- (1.302,1.982);
      \node[left=4pt,defscol,font=\small\bfseries,fill=white,draw=defscol,
            line width=0.3pt,inner sep=2pt,rounded corners=1pt]
        at (1.555,1.357) {$\vec{N}$};
      \draw[-{Stealth[length=11pt,width=7pt]}, line width=2.5pt, black!65]
        (1.714,0.963) -- (1.714,-0.22);
      \node[right=4pt,black!65,font=\small\bfseries,fill=white,draw=black!50,
            line width=0.3pt,inner sep=2pt,rounded corners=1pt]
        at (1.714,0.37) {$\vec{G}$};
      \draw[-{Stealth[length=11pt,width=7pt]}, line width=2.5pt, red!70!black]
        (1.808,0.731) -- (0.974,0.393);
      \node[below=5pt,red!70!black,font=\small\bfseries,fill=white,draw=red!70!black,
            line width=0.3pt,inner sep=2pt,rounded corners=1pt]
        at (1.391,0.562) {$\vec{R}$};
      \draw[dashed,black!45,semithick] (1.808,0.731) -- (1.808,1.982);
      \draw[dashed,black!45,semithick] (1.808,1.982) -- (1.302,1.982);
      \node[right=2pt,black!55,font=\scriptsize] at (1.808,1.357) {$N_V$};
      \node[above=2pt,black!55,font=\scriptsize] at (1.555,1.982) {$N_H$};
      \draw[-{Stealth[length=8pt,width=5pt]},black!70,semithick]
        (3.2,0.08) -- (3.2,0.95);
      \node[right=2pt,black!70,font=\scriptsize] at (3.2,0.95) {$V$};
      \draw[-{Stealth[length=8pt,width=5pt]},black!70,semithick]
        (3.2,0.08) -- (2.40,0.08);
      \node[below=2pt,black!70,font=\scriptsize] at (2.40,0.08) {$H$};
      \fill[black!70] (3.2,0.08) circle (1.5pt);
      \fill[defscol] (1.808,0.731) circle (2pt);
      \fill[black!55] (1.714,0.963) circle (2pt);
    \end{tikzpicture}
      \caption{\footnotesize Dosert veibane i tverrsnitt. $\vec{N}$ dekomponeres i $N_V$ (opp) og $N_H$ (inn mot sentrum). Friksjonskraften $\vec{R}$ peker nedover veibanen.}
      \end{figure}
    \end{column}
  \end{columns}
\end{frame}

%--------------------------------------------------------------------------
%  SLIDE 10 — Oppsummering
%--------------------------------------------------------------------------
\begin{frame}{}

  \oppsum{Krefter i sirkelbevegelse}{%
  \begin{enumerate}[1.]
    \item[] Sentripetalkraft: $\sum F = m v^2/r$, rettet inn mot sentrum.
    \item[] Dosert veibane (uten friksjon): $\tan\alpha = v^2/(r\,\mathrm{g})$.
  \end{enumerate}
  }

\end{frame}

%--------------------------------------------------------------------------
%  SLIDE 11 — Neste gang
%--------------------------------------------------------------------------
\begin{frame}
\begin{center}
  \Huge{Neste gang: Bevaring av bevegelsesmengde}
  \begin{itemize}
    \item Isolerte systemer
    \item Rekyl og støt
  \end{itemize}
  \vspace{6mm}
  {\normalsize\itshape Det som går inn, må komme ut\ldots}
\end{center}
\end{frame}

\end{document}
